\documentclass[12pt,a4paper]{report}

% ── Packages ──────────────────────────────────────────────────────────────────
\usepackage[margin=1in]{geometry}
\usepackage{graphicx}
\usepackage{booktabs}
\usepackage{longtable}
\usepackage{array}
\usepackage{tabularx}
\usepackage{setspace}
\usepackage{parskip}
\usepackage{titlesec}
\usepackage{fancyhdr}
\usepackage{enumitem}
\usepackage{amsmath}
\usepackage{microtype}
\usepackage[T1]{fontenc}
\usepackage{lmodern}
\usepackage{xcolor}
\usepackage{caption}
\usepackage{tocloft} % Added to customize Table of Contents structure

% ── TOC Customization (Fixes Screenshot 2026-06-04 072015.png layout mismatch) 
\renewcommand{\cftchapleader}{\cftdotfill{\cftdotsep}} % Forces dot leaders for all chapters in TOC

% ── Package Configuration & Link Colors ─────────────────────────────────────
\usepackage{hyperref}
\hypersetup{
    colorlinks=true,
    linkcolor=black,     % Set to black for professional report printing
    urlcolor=blue
}

% ── Page style & Layout Fixes ─────────────────────────────────────────────────
\pagestyle{fancy}
\fancyhf{}
\setlength{\headheight}{14.5pt} 
\addtolength{\topmargin}{-2.5pt} 

\rhead{\small AI Virtual Mouse}
\lhead{\small Dr. Bapuji Salunkhe Institute of Engineering and Technology}
\cfoot{\thepage}

% ── Section formatting ────────────────────────────────────────────────────────
\titleformat{\chapter}[display]
  {\normalfont\large\bfseries\centering}{Chapter \thechapter}{12pt}{\Large}
\titlespacing*{\chapter}{0pt}{20pt}{20pt}

\titleformat{\section}{\normalfont\large\bfseries}{\thesection}{1em}{}
\titleformat{\subsection}{\normalfont\normalsize\bfseries}{\thesubsection}{1em}{}

\onehalfspacing

% ══════════════════════════════════════════════════════════════════════════════
\begin{document}

% ── Title Page ────────────────────────────────────────────────────────────────
\begin{titlepage}
  \centering
  \vspace*{1cm}

  {\Large\bfseries A mini project report on}\\[0.4cm]
  {\LARGE\bfseries ``AI VIRTUAL MOUSE''}\\[0.6cm]

  {\normalsize Submitted in Partial Fulfillment for the Requirement of}\\[0.2cm]
  {\large\bfseries Third Year}\\[0.1cm]
  {\normalsize In}\\[0.1cm]
  {\large\bfseries Computer Science and Engineering}\\[0.6cm]

  {\normalsize Submitted to}\\[0.4cm]
  {\large\bfseries Dr.\ Babasaheb Ambedkar Technological University, Lonere}\\[0.6cm]

  {\normalsize By}\\[0.4cm]
  \begin{tabular}{ll}
    Pranav Gavali       & (23064681242046) \\
    Varad Sawant        & (23064681242048) \\
    Bhumika Dhumal      & (24064681242505) \\
    Samruddhi Majgaonkar & (24064681242504) \\
    Prachi Shinde       & (24064681242503) \\
  \end{tabular}\\[0.6cm]

  {\normalsize Under the Guidance of}\\[0.2cm]
  {\large\bfseries Prof.\ V.S.\ Jadhav}\\[0.8cm]

  {\large\bfseries Department of Computer Science and Engineering}\\[0.2cm]
  {\large\bfseries Dr.\ Bapuji Salunkhe Institute of Engineering and Technology,}\\
  {\normalsize 2130, E ward, Tarabai Park, Kolhapur}\\[0.4cm]
  {\large\bfseries Academic Year 2025--2026}

  \vfill
  \hrule\vspace{0.3cm}
  {\small External Examiner: \underline{\hspace{4cm}} \hfill Department of \underline{\hspace{4cm}}}
\end{titlepage}

% ── Certificate ───────────────────────────────────────────────────────────────
\chapter*{Certificate}
\addcontentsline{toc}{chapter}{Certificate}
\thispagestyle{plain}

\begin{center}
  {\large\bfseries Department of Computer Science and Engineering}\\[0.2cm]
  {\normalsize Shree Swami Vivekanand Shikshan Sanstha's\\
  Dr.\ Bapuji Salunkhe Institute of Engineering and Technology}\\[0.6cm]
  {\LARGE\bfseries\color{red} CERTIFICATE}
\end{center}

\vspace{0.5cm}
This is to certify that the mini project entitled:

\begin{center}
  {\large\bfseries\color{red} ``AI VIRTUAL MOUSE''}
\end{center}

Submitted by the following students:

\begin{center}
  \begin{tabular}{ll}
    Pranav Gavali        & (23064681242046) \\
    Varad Sawant         & (23064681242048) \\
    Bhumika Dhumal       & (24064681242505) \\
    Samruddhi Majgaonkar & (24064681242504) \\
    Prachi Shinde        & (24064681242503) \\
  \end{tabular}
\end{center}

has been successfully completed under the guidance of \textbf{Prof.\ V.S.\ Jadhav} during the
\textbf{Academic Year 2025--2026}. This report is submitted in partial fulfillment of the
requirements for the Third Year Mini Project in Computer Science and Engineering.

\vspace{1.5cm}
\begin{tabular}{p{6cm}p{6cm}}
  \textbf{Dr.\ R.\ D.\ Bhosale}      & \textbf{Prof.\ V.S.\ Jadhav} \\
  Project Coordinator                 & Project Guide                 \\
\end{tabular}

\vspace{1.5cm}
\begin{tabular}{p{5cm}p{4cm}p{4cm}}
  \textbf{Mr.\ A.\ R.\ Sawant}        & \textbf{Dr.\ S.\ G.\ Sapate} & \\
  H.O.D.                              & Principal                     & External Examiner \\
  Dept.\ of Computer Science \& Engg. & & \\
\end{tabular}

% ── Acknowledgement ───────────────────────────────────────────────────────────
\chapter*{Acknowledgement}
\addcontentsline{toc}{chapter}{Acknowledgement}

We take this opportunity to express our sincere gratitude to our respected project guide,
\textbf{Prof.\ V.\ S.\ Jadhav}, for providing us with the chance to work on this image processing
project and for offering feedback and suggestions that helped make this project feasible. This
project could not have been accomplished without her leadership. Her insightful feedback,
constructive criticism, and timely suggestions were instrumental in shaping the direction of
this research.

We express our sincere thanks to the respected Principal, \textbf{Dr.\ S.G.\ Sapate}, for
providing the necessary facilities and academic environment.

We are thankful to the respected Head of Department, \textbf{Prof.\ A.R.\ Sawant}, Department
of Computer Engineering, for departmental support and encouragement.

We are also thankful to all teaching and non-teaching staff members of the department for their
cooperation and support.

Lastly, we want to express our gratitude to our close companions and family members for their
inspiration and steadfast support throughout this journey.

\vspace{1cm}
\begin{flushright}
  Pranav P.\ Gavali (43139)\\
  Varad N.\ Sawant (43142)\\
  Prachi P.\ Shinde (43143)\\
  Samruddhi S.\ Majgaonkar (43144)\\
  Bhumika S.\ Dhumal (43146)
\end{flushright}

% ── Abstract ──────────────────────────────────────────────────────────────────
\chapter*{Abstract}
\addcontentsline{toc}{chapter}{Abstract}

The AI Virtual Mouse is a gesture-based controller application that utilizes hand gestures
captured through a webcam to perform various system control functions, leveraging the
\textbf{MediaPipe} library for hand tracking and gesture recognition. The application is
structured into \textit{Gesture Encodings}, \textit{Hand Recognition}, and \textit{Controller}
classes, and interprets hand movements to execute corresponding commands such as cursor
movement, clicking, scrolling, and system parameter adjustments including volume and
brightness control.

The system discerns a range of hand gestures including fist, V-gesture, and pinch to enable
diverse actions. It intelligently differentiates between major and minor hands, assigning
specific functionalities accordingly. The \texttt{GestureController} class manages webcam
initialization, hand landmark capture, and processing. This application presents a real-time
interface for intuitive system interaction via hand gestures, offering a novel means of
navigating and controlling computer functions without the need for physical input devices.

\newpage  % ← forces Table of Contents onto a new page

% ── Table of Contents ─────────────────────────────────────────────────────────
\tableofcontents

% ══════════════════════════════════════════════════════════════════════════════
\chapter{Introduction}

In a world driven by technological innovation, the way we interact with our digital devices
plays a pivotal role in our daily lives. Traditional input devices, such as mice and touchpads,
have long been the standard for navigating the digital realm. However, as the frontiers of
technology continue to expand, a new paradigm in human-computer interaction emerges in the
form of the AI Virtual Mouse.

The AI Virtual Mouse represents a groundbreaking fusion of artificial intelligence (AI),
computer vision, and machine learning, designed to revolutionize the way we control and
navigate digital interfaces. This cutting-edge technology eliminates the need for physical
pointing devices and, instead, leverages the subtlety of our hand movements and gestures to
manipulate the digital world.

\section{Background of Concept}

The concept behind the AI Virtual Mouse is to create a gesture-based controller using
computer vision techniques. The system utilizes the MediaPipe library to detect and track
hand movements in real-time using a webcam. By analyzing the hand landmarks detected by the
library, the system maps specific hand gestures to predefined actions or commands, such as
controlling the mouse cursor, clicking, scrolling, adjusting system settings like brightness
and volume, and more. The system provides visual feedback by overlaying hand landmarks and
recognized gestures on the webcam feed, allowing users to monitor their interactions. Overall,
the system demonstrates an intuitive and hands-free way for users to interact with their
computers through gestures, enhancing accessibility and user experience.

\section{Problem Specification}

The problem addressed in this project is the need for a more accessible and efficient cursor
control system through the development and implementation of an AI-driven virtual mouse.
Traditional input methods such as physical mice and touchpads are impractical or limiting for
certain users, especially those with disabilities.

The system should accurately recognize hand gestures in real-time and map them to specific
actions or commands, such as controlling the mouse cursor, clicking, scrolling, adjusting
system settings like brightness and volume, and more. The system should provide a
user-friendly interface and ensure smooth and intuitive interaction. Additionally, the system
should be robust to noise and variations in hand movements, providing reliable performance
across different environments and conditions.

% ══════════════════════════════════════════════════════════════════════════════
\chapter{Related Studies / Literature Survey}

\section{Literature Survey}

Several researchers and developers have worked on gesture-based computer interaction systems.
Numerous studies have investigated the use of computer vision techniques, such as hand
tracking and landmark detection, to interpret human gestures for Human-Computer Interaction
(HCI) applications.

\begin{longtable}{>{\bfseries}p{3cm} p{3.5cm} p{2.5cm} p{4.5cm}}
\caption{Summary of Related Work} \label{tab:litsurvey} \\
\toprule
\textbf{Author} & \textbf{System / Method} & \textbf{Technology} & \textbf{Outcome / Limitation} \\
\midrule
\endfirsthead
\toprule
\textbf{Author} & \textbf{System / Method} & \textbf{Technology} & \textbf{Outcome / Limitation} \\
\midrule
\endhead
\bottomrule
\endfoot
Cao et al.\ (2018)  & OpenPose -- real-time multi-person body/hand keypoint detection & Deep Learning, Computer Vision & Enabled gesture recognition for various tasks; computationally intensive. \\[4pt]
Simon et al.\ (2017) & Real-time hand gesture recognition using depth data & Depth Sensors, CNN & Facilitated intuitive interactions in virtual environments; required special hardware. \\[4pt]
Sun et al.\ (2019)  & Deep learning-based hand gesture recognition method & Deep Learning & Achieved high performance in real-world scenarios. \\[4pt]
Zhang et al.\ (2020) & CNN-based hand gesture recognition & Convolutional Neural Networks & Superior performance compared to traditional approaches. \\[4pt]
Guna et al.\ (2014) & Gesture-based interaction in healthcare settings & Computer Vision & Enabled hands-free control of medical devices for healthcare professionals. \\
\end{longtable}

\section{Existing Systems}

\begin{longtable}{p{3.5cm} p{4cm} p{5cm}}
\caption{Existing Systems Comparison} \label{tab:existing} \\
\toprule
\textbf{System / Product} & \textbf{Main Feature} & \textbf{Limitation} \\
\midrule
\endfirsthead
\toprule
\textbf{System / Product} & \textbf{Main Feature} & \textbf{Limitation} \\
\midrule
\endhead
\bottomrule
\endfoot
Physical Mouse & Standard cursor control & Not accessible for users with mobility impairments. \\[4pt]
Touchpad / Touchscreen & Touch-based input & Requires physical contact; limited gesture variety. \\[4pt]
Leap Motion & Depth-based hand tracking & Requires specialized hardware; high cost. \\[4pt]
Voice Control (Siri/Cortana) & Voice-based system control & Limited cursor control; noisy environment dependency. \\
\end{longtable}

\section{Research Gap}

\begin{itemize}[noitemsep]
  \item Existing systems often require specialized hardware (e.g., depth cameras).
  \item Many systems are not easily deployable using only a standard webcam.
  \item Approval and accessibility workflows for institutional use are rarely addressed.
  \item Real-time gesture recognition with volume and brightness control in a single system is uncommon.
\end{itemize}

\section{Problem Statement}

To design and implement a webcam-based AI Virtual Mouse system that uses hand gesture recognition
to control computer cursor movement, mouse clicks, scrolling, and system settings (volume and
brightness) in real-time, using only standard computer hardware, without requiring any specialized
physical input devices.

\section{Specific Objectives}

\begin{enumerate}[noitemsep]
  \item To implement real-time hand tracking using the MediaPipe library.
  \item To develop gesture encoding and recognition for various hand gestures.
  \item To map recognized gestures to system control actions (cursor movement, clicking, scrolling).
  \item To implement system-level controls including volume and brightness adjustment.
  \item To differentiate between dominant and non-dominant hands for adaptive control.
  \item To provide real-time visual feedback via webcam overlay.
\end{enumerate}

\section{Scope of Work}

The proposed system focuses on webcam-based hand gesture recognition for computer control. It
includes gesture detection, cursor control, click actions, scroll actions, volume and brightness
control, and real-time visual feedback. The system is designed for Windows OS and requires a
standard webcam.

\section{Limitations}

\begin{itemize}[noitemsep]
  \item Mobile application support is not implemented.
  \item System performance may vary depending on lighting conditions.
  \item Background clutter may reduce gesture recognition accuracy.
  \item The gesture vocabulary is limited to predefined gestures only.
  \item Continuous use may cause hand fatigue.
\end{itemize}

% ══════════════════════════════════════════════════════════════════════════════
\chapter{Proposed Methodology and Design}

\section{System Overview}

The proposed AI Virtual Mouse system uses a webcam-based client architecture. Users perform hand
gestures in front of the webcam. The system captures video frames, processes them using the
MediaPipe library to detect hand landmarks, recognizes gestures using the \texttt{HandRecog}
class, and executes corresponding computer control commands via the \texttt{Controller} class.
The system provides real-time visual feedback by overlaying detected landmarks on the live
webcam feed.

\section{System Architecture / Block Diagram}

The system is organized into three main layers:
\begin{itemize}[noitemsep]
  \item \textbf{Video Capture} -- input via webcam
  \item \textbf{Gesture Processing} -- MediaPipe, HandRecog
  \item \textbf{Action Execution} -- Controller class
\end{itemize}

The \texttt{GestureController} class orchestrates the entire pipeline from frame capture to
command execution.

\section{Data Flow Diagrams}

\subsection{DFD Level 0}

At level 0, the system takes \textit{User Movements} as input, performs \textit{Motion
Analysis} as processing, and produces \textit{Cursor Movements} as output. The entire AI
Virtual Mouse system is shown as a single process.

\subsection{DFD Level 1}

At level 1, the DFD expands to show the internal processes: the User starts the Gesture
Controller, which checks for a valid webcam connection, captures video frames, detects hand
gestures, recognizes them, and triggers corresponding system actions.

\subsection{DFD Level 2}

At level 2, the DFD further details the gesture recognition loop: the Gesture Controller
interacts with the Webcam to capture frames, passes them to the Gesture Recognition module,
processes recognized Hand Gestures, triggers actions (Execute Commands), and can release
camera resources to terminate the process.

\section{Class Diagram}

The system consists of three main classes:

\begin{itemize}[noitemsep]
  \item \textbf{GestureController} -- manages webcam, frame processing, and hand classification
  \item \textbf{HandRecog} -- processes landmarks and recognizes gestures
  \item \textbf{Controller} -- executes system actions based on gestures
\end{itemize}

\texttt{GestureController} uses both \texttt{HandRecog} and \texttt{Controller}.

\section{UML Diagram}

The UML diagram shows the \texttt{GestureController} as the central component with methods for
detecting, encoding, and handling gestures, and performing actions. It is associated with
\texttt{MediaPipeLibrary}, \texttt{Gest} (gesture encodings), \texttt{User}, and
\texttt{Computer} entities.

\section{Sequence Diagram}

The sequence diagram shows the interaction flow:

\begin{enumerate}[noitemsep]
  \item User starts the Python script
  \item System initializes
  \item Webcam initializes
  \item Loop begins
  \item System captures video frames
  \item MediaPipe processes frames and detects hand gestures
  \item System classifies and recognizes gestures
  \item PyAutoGUI triggers corresponding actions
  \item User receives control over computer functions
\end{enumerate}

\section{Algorithm of Proposed System}

\begin{enumerate}
  \item \textbf{Initialization:} Initialize the \texttt{GestureController} by setting up the
    camera and necessary libraries such as OpenCV, MediaPipe, PyAutoGUI, etc.

  \item \textbf{Hand Detection and Tracking:} Continuously capture frames from the webcam and
    use the MediaPipe library to detect and track hand landmarks in each frame. Identify the
    dominant hand based on the classification provided by MediaPipe.

  \item \textbf{Gesture Recognition:} Process the hand landmarks to recognize various hand
    gestures, such as fist, palm, pinch, finger gestures, etc. Map these gestures to
    predefined commands or actions.

  \item \textbf{Gesture Control:} Based on the recognized gestures, execute corresponding
    actions on the computer. Actions may include moving the mouse cursor, clicking, scrolling,
    adjusting system settings (e.g., brightness, volume), etc.

  \item \textbf{Feedback and Visualization:} Provide visual feedback to the user by overlaying
    hand landmarks and recognized gestures on the webcam feed.

  \item \textbf{User Interaction:} Allow the user to interact with the system by performing
    gestures, observing the feedback on the screen, and controlling the computer accordingly.

  \item \textbf{Termination:} Terminate the gesture controller when the user decides to exit
    or press the Enter key to close the application.
\end{enumerate}

% ══════════════════════════════════════════════════════════════════════════════
\chapter{Experimental Setup and Implementation Details}

\section{Software Requirements and Specifications}

\begin{table}[ht]
\centering
\caption{Software Requirements}
\begin{tabular}{ll}
\toprule
\textbf{Component} & \textbf{Specification} \\
\midrule
Operating System       & Windows 11 \\
Programming Language   & Python \\
Computer Vision Library & OpenCV (cv2), MediaPipe \\
Mouse/Keyboard Control & PyAutoGUI \\
Volume Control         & pycaw, comtypes, ctypes \\
Brightness Control     & screen\_brightness\_control \\
Data Encoding          & Google Protocol Buffers \\
Math Library           & Python math module \\
Enum Library           & Python Enum (IntEnum) \\
\bottomrule
\end{tabular}
\end{table}

\section{Hardware Requirements}

\begin{table}[ht]
\centering
\caption{Hardware Requirements}
\begin{tabular}{ll}
\toprule
\textbf{Component} & \textbf{Specification} \\
\midrule
Processor & AMD Ryzen 5 5600H with Radeon Graphics 3.30\,GHz \\
RAM       & 8.00\,GB \\
Webcam    & Standard built-in or external USB webcam \\
Keyboard  & Standard Windows Keyboard \\
Mouse     & 2 or 3 Button Mouse (for fallback) \\
\bottomrule
\end{tabular}
\end{table}

\section{Technology Details}

\begin{description}[noitemsep, leftmargin=2cm]
  \item[\textbf{OpenCV (cv2)}] Used for real-time image processing tasks like capturing video
    frames and image manipulation.
  \item[\textbf{MediaPipe}] A machine learning framework by Google used for hand tracking and
    landmark detection.
  \item[\textbf{PyAutoGUI}] A Python library for cross-platform mouse and keyboard simulation
    based on detected gestures.
  \item[\textbf{Math Module}] Used for mathematical calculations such as computing distances
    between landmarks.
  \item[\textbf{ctypes / comtypes}] Used for interacting with Windows Core Audio APIs to
    control system volume.
  \item[\textbf{pycaw}] Python library for interacting with the Windows Core Audio APIs for
    volume control.
  \item[\textbf{screen\_brightness\_control}] Used for programmatically adjusting screen
    brightness.
  \item[\textbf{Google Protobuf}] Used for converting hand tracking results from MediaPipe
    into a dictionary format.
\end{description}

\section{Project Plan}

\begin{table}[ht]
\centering
\caption{Project Plan}
\begin{tabularx}{\textwidth}{lX}
\toprule
\textbf{Work Task} & \textbf{Description} \\
\midrule
Literature Research & Related work done for research about needs of customers \\
System Analysis     & Critical analysis and comparison of technologies and results achieved in research \\
Design \& Planning  & Modeling, design, and dataset searching/creation \\
Implementation      & Implementation of the literature review and proposed methodology \\
Initial Report      & Prepare and upload initial report \\
Final Report        & Prepare and upload final report \\
\bottomrule
\end{tabularx}
\end{table}

\section{Implementation Details}

The implementation includes: camera initialization using OpenCV, hand landmark detection using
MediaPipe Hands module, gesture classification using the \texttt{HandRecog} class, execution
of system commands (cursor movement, clicks, scroll, volume, brightness) using the
\texttt{Controller} class, and real-time visual feedback via the OpenCV display window titled
\textit{Gesture Controller}.

The \texttt{GestureController} class manages the main loop: it captures frames, flips them
horizontally (mirror effect), converts to RGB for MediaPipe, processes results, classifies
hands as major/minor, and calls the Controller to handle the corresponding actions.

% ══════════════════════════════════════════════════════════════════════════════
\chapter{Results Analysis and Discussion}

\section{Test Specification}

\subsection{Whitebox Testing}

Unit Testing, Integration Testing, Boundary Testing, and Error Handling Testing were performed
on individual components including gesture recognition, hand tracking, and controller actions.

\begin{itemize}[noitemsep]
  \item \textbf{Unit Testing:} Test individual components such as gesture recognition, hand
    tracking, and controller actions in isolation to ensure they function correctly.
  \item \textbf{Integration Testing:} Verify the interaction between different modules to
    ensure they integrate seamlessly.
  \item \textbf{Boundary Testing:} Test with extreme inputs or boundary conditions to validate
    behavior under these scenarios.
  \item \textbf{Error Handling Testing:} Evaluate how the system responds to unexpected inputs
    or errors.
\end{itemize}

\subsection{Blackbox Testing}

Functionality, Boundary, Error Handling, and Compatibility testing were performed from the
user perspective.

\begin{itemize}[noitemsep]
  \item \textbf{Functionality Testing:} Ensured the script accurately detects various hand
    gestures and each recognized gesture triggers the corresponding action.
  \item \textbf{Boundary Testing:} Tested extreme cases such as detecting gestures with very
    low or high hand visibility and at the edge of the webcam frame.
  \item \textbf{Compatibility Testing:} Tested the script on different operating systems and
    verified compatibility with various webcam devices.
\end{itemize}

\section{Test Cases}

\begin{longtable}{c p{4cm} p{3cm} p{3cm} c}
\caption{Test Cases} \label{tab:testcases} \\
\toprule
\textbf{TC ID} & \textbf{Test Case} & \textbf{Expected Result} & \textbf{Actual Output} & \textbf{Status} \\
\midrule
\endfirsthead
\toprule
\textbf{TC ID} & \textbf{Test Case} & \textbf{Expected Result} & \textbf{Actual Output} & \textbf{Status} \\
\midrule
\endhead
\bottomrule
\endfoot
1  & Camera initializes correctly              & Initializes without errors       & Initialized without errors      & Pass \\
2  & Camera captures video frames              & Frames captured successfully     & Frames captured successfully    & Pass \\
3  & Frames are flipped horizontally           & Frames flipped horizontally      & Flipped as expected             & Pass \\
5  & Hand landmarks are detected               & Landmarks detected accurately    & Detected accurately             & Pass \\
6  & Single hand detection                     & Only one hand detected           & One hand detected as expected   & Pass \\
7  & Multiple hands detected                   & Multiple hands detected          & Detected as expected             & Pass \\
9  & Gestures correctly encoded                & Gestures encoded accurately      & Encoded accurately              & Pass \\
12 & Mouse movements smooth/responsive         & Smooth and responsive            & Smooth and responsive           & Pass \\
13 & Left-click functionality                  & Left-click performed accurately  & Left-click performed accurately & Pass \\
14 & Right-click functionality                 & Right-click performed accurately & Right-click performed accurately & Pass \\
15 & Double-click functionality                & Double-click performed accurately & Double-click performed accurately & Pass \\
19 & Brightness control works                  & Brightness adjusts as expected   & Brightness adjusts as expected  & Pass \\
21 & Graceful termination on exit              & System terminates gracefully     & Terminates gracefully           & Pass \\
22 & Compatibility across OS                   & Functions correctly on different OS & Functions correctly           & Pass \\
25 & Users of different ages                   & Functions correctly with all ages & Functions correctly            & Pass \\
\end{longtable}

\section{Statistical Analysis}

The system was tested with sample gesture data for various operations. The success rate is
defined as:

\[
  \text{Success Rate} = \frac{\text{Successful Operations}}{\text{Total Operations}} \times 100
\]

\begin{table}[ht]
\centering
\caption{Statistical Analysis of System Performance}
\begin{tabular}{lrrr}
\toprule
\textbf{Operation} & \textbf{Total Tests} & \textbf{Successful Tests} & \textbf{Success Rate} \\
\midrule
Gesture Detection          & 50 & 49 & 98\% \\
Cursor Movement            & 50 & 50 & 100\% \\
Left Click                 & 50 & 49 & 98\% \\
Right Click                & 50 & 48 & 96\% \\
Volume/Brightness Control  & 50 & 49 & 98\% \\
Scroll Action              & 50 & 48 & 96\% \\
\bottomrule
\end{tabular}
\end{table}

\section{Comparative Analysis}

\begin{table}[ht]
\centering
\caption{Comparative Analysis of Input Systems}
\begin{tabularx}{\textwidth}{Xcccc}
\toprule
\textbf{System} & \textbf{Physical Device} & \textbf{Accessibility} & \textbf{Gesture Support} & \textbf{Real-time Feedback} \\
\midrule
Physical Mouse         & Yes            & Limited   & None              & No  \\
Touchpad               & Yes            & Limited   & Basic             & No  \\
Leap Motion            & Yes (specialized) & Good   & Advanced          & Yes \\
\textbf{Proposed System} & Webcam only  & Excellent & Multiple gestures & Yes \\
\bottomrule
\end{tabularx}
\end{table}

\section{Discussion on Successful and Failure Cases}

\textbf{Successful cases} include: valid hand detection in well-lit environments, correct
gesture recognition for cursor movement and clicks, smooth cursor control, and accurate
volume/brightness adjustments.

\textbf{Failure cases} may occur due to: poor lighting conditions causing low detection
accuracy, hand occlusion or unusual angles, background clutter interfering with landmark
detection, and minor inaccuracies in right-click and drag-and-drop operations. These can be
mitigated by improving the fingertip detection algorithm and adding adaptive lighting
compensation.

% ──══════════════════════════════════════════════════════════════════════════
\chapter{Conclusion and Future Scope}

\section{Advantages}

\begin{itemize}[noitemsep]
  \item \textbf{Hands-Free Interaction:} Allows users to control their computer without any
    physical input devices.
  \item \textbf{Intuitive Interface:} Gesture-based controls provide a natural and intuitive
    way to interact with the computer.
  \item \textbf{Real-Time Feedback:} Provides immediate feedback to users based on their hand
    gestures, enhancing user experience.
  \item \textbf{Flexible Applications:} Can be adapted for various tasks such as cursor
    movement, clicking, scrolling, and system control.
  \item \textbf{Accessibility:} Offers an alternative input method for individuals with
    mobility impairments.
\end{itemize}

\section{Disadvantages}

\begin{itemize}[noitemsep]
  \item \textbf{Learning Curve:} Users may need time to learn and memorize the supported hand
    gestures and their corresponding actions.
  \item \textbf{Accuracy Limitations:} Detection accuracy may vary depending on lighting
    conditions, background clutter, and hand positioning.
  \item \textbf{Fatigue:} Continuous gesture-based interaction may cause hand fatigue over
    time.
  \item \textbf{Limited Gesture Vocabulary:} The number of distinct gestures that can be
    reliably recognized may be limited.
  \item \textbf{Dependency on Hardware:} Relies on the availability and quality of a
    compatible webcam.
\end{itemize}

\section{Future Scope}

\begin{itemize}[noitemsep]
  \item \textbf{Voice Control Integration:} Combine gesture and voice recognition to allow
    multimodal computer control.
  \item \textbf{Improved Accessibility:} Continue improvements for users with disabilities as
    AI technology advances.
  \item \textbf{Mobile Application:} Develop a mobile version of the gesture controller for
    smartphones and tablets.
  \item \textbf{Facial Recognition Integration:} Combine with facial recognition for
    multi-input control.
  \item \textbf{Cloud Deployment:} Enable cloud-based gesture recognition for wider
    accessibility.
  \item \textbf{Expanded Gesture Vocabulary:} Improve the algorithm to support a wider set of
    reliable gestures.
\end{itemize}

\section{Conclusion}

This mini project presented an AI Virtual Mouse system using hand gesture recognition for
computer control. The system provides real-time hand tracking, gesture recognition, cursor
movement, mouse clicks (left, right, double), scroll actions, and volume and brightness
control, all using only a standard webcam.

The proposed system eliminates the need for a physical mouse and improves accessibility.
Using MediaPipe for hand detection and tracking and PyAutoGUI for system control, we
successfully implemented a real-time gesture-based computer control system. The system has
some limitations such as minor accuracy reduction in right-click and drag-and-drop functions;
these will be addressed in future work by improving the fingertip detection algorithm to
produce more accurate results.

% ── References ────────────────────────────────────────────────────────────────
\begin{thebibliography}{9}

\bibitem{cao2021}
Z.~Cao, G.~Hidalgo, T.~Simon, S.-E.~Wei, and Y.~Sheikh,
``OpenPose: Realtime Multi-Person 2D Pose Estimation Using Part Affinity Fields,''
\textit{IEEE Transactions on Pattern Analysis and Machine Intelligence}, vol.\ 43, no.\ 1, pp.\ 172--186, Jan.\ 2021.

\bibitem{simon2017}
T.~Simon, H.~Joo, I.~Matthews, and Y.~Sheikh,
``Hand Keypoint Detection in Single Images using Multiview Bootstrapping,''
\textit{Proceedings of the IEEE Conference on Computer Vision and Pattern Recognition (CVPR)}, pp.\ 1145--1153, 2017.

\end{thebibliography}

\end{document}
